% ============================================================
% ΔΙΟΡΘΩΜΕΝΗ έκδοση — Κεφάλαιο 0
% Μιγαδικοί Αριθμοί
%
% Αφαιρέθηκαν τα \begin{ekfonisicas} / \begin{ekfonisipy}
% ώστε να μην εμφανίζεται κάθετη γραμμή μέσα στα box.
%
% QR code: qr_cf0b663610cc.png  (→ latex/qrcodes/)
% URL: https://nikmatz.github.io/ELADIC_GR/code/ch00/
% ============================================================

\casactivbox%
  {Maxima: Μιγαδικοί — Εκθετική Μορφή, Φάσορες, De Moivre}%
  {%
    \label{cas:ch0_complex_exponential}
    Στην ανάλυση AC κυκλωμάτων χρησιμοποιούμε \textbf{φάσορες} (phasors)
    σε εκθετική μορφή $z = r\,e^{i\theta}$.
    Δίνονται $z_A = 5\,e^{i\pi/3}$ και $z_B = 2\,e^{i\pi/6}$.
    \begin{enumerate}[label=(\alph*),itemsep=2pt]
      \item Μετατρέψτε $z_A$, $z_B$ σε ορθογώνια μορφή $a+bi$
        (\texttt{rectform}).
      \item Υπολογίστε $z_A \cdot z_B$ και $z_A / z_B$. Επαληθεύστε
        πρόσθεση/αφαίρεση γωνιών και πολλαπλασιασμό/διαίρεση μέτρων.
      \item Επαληθεύστε De Moivre για $n=3$:
        $\bigl(5\,e^{i\pi/3}\bigr)^3 = -125$.
      \item Βρείτε τις τρεις ρίζες $x^3 = 8$ (\texttt{solve}).
        Σχολιάστε τη γεωμετρική τους διάταξη στο επίπεδο.
    \end{enumerate}
    \smallskip
    \textbf{Εντολές Maxima:}
    \texttt{rectform(z)}, \texttt{abs(z)}, \texttt{carg(z)},
    \texttt{simplify(\ldots)}, \texttt{solve(x\^{}3=8, x)}
  }%
  {https://nikmatz.github.io/ELADIC_GR/code/ch00/}%
  {qr_cf0b663610cc.png}

\subsection*{Δραστηριότητα Python}

\begin{pybox}{Python: Μιγαδικοί — cmath, SymPy, Διάγραμμα Argand}%
             {https://nikmatz.github.io/ELADIC_GR/code/ch00/}%
             {qr_cf0b663610cc.png}
\label{py:ch0_complex}
Δίνονται $z_1 = 3+4i$, $z_2 = 1-2i$ και φάσορες
$z_A = 5\,e^{i\pi/3}$, $z_B = 2\,e^{i\pi/6}$.
\begin{enumerate}[label=(\alph*),itemsep=2pt]
  \item Υπολογίστε βασικές πράξεις με \texttt{cmath}:
    $z_1+z_2$, $z_1 \cdot z_2$, $|z_1|$, $\arg(z_1)$.
  \item Για τους φάσορες: υπολογίστε γινόμενο και πηλίκο,
    επαληθεύστε νόμους μέτρου και γωνίας.
  \item Βρείτε τις δύο τετραγωνικές ρίζες του $1+i$
    και επαληθεύστε $x^2 = 1+i$.
  \item Επαληθεύστε De Moivre $n=3$ με SymPy (συμβολικά).
  \item \textit{Προαιρετικό:} Σχεδιάστε διάγραμμα Argand
    με τις πράξεις και τις ρίζες $x^3=8$ (Matplotlib).
\end{enumerate}
\medskip
\textbf{Βιβλιοθήκες / Εντολές:}
\texttt{cmath.polar(z)}, \texttt{cmath.rect(r,$\theta$)},
\texttt{cmath.phase(z)}, \texttt{sympy.solve()},
\texttt{matplotlib.pyplot}
\end{pybox}
